\bookchapter{第四十三章：对话工程——与AI高效协作的对话技术}

\chapterintro{对话是 AI 编程的核心交互方式。但请注意：“对话”不等于“聊天”——高效的 AI 编程对话是一种有结构的沟通技术，我们称之为“对话工程”。它研究的是：如何用最少的对话轮次，获得最高质量的代码产出。说白了，这一章就是教你怎么和 AI“把话说清楚”。}

\booksection{一、Prompt 分层策略}

与 AI 的对话可以分为三个层次，每层承担不同的功能。理解这三个层次，你就能像搭积木一样组合出高效的 Prompt。

\begin{center}
\begin{tabularx}{\linewidth}{|X|X|X|X|}
\hline
\textbf{层次} & \textbf{功能} & \textbf{示例} & \textbf{频率} \tabularnewline
\hline
系统层 & 设定AI的角色和全局约束 & "你是一个React专家，代码必须用TypeScript" & 项目开始时设置一次 \tabularnewline
任务层 & 描述当前要做的具体任务 & "实现一个表单组件，包含用户名和密码字段" & 每个任务开始时给出 \tabularnewline
反馈层 & 对AI产出进行评价和修正 & "密码字段缺少显示/隐藏切换，请添加" & 根据产出质量给出 \tabularnewline
\hline
\end{tabularx}
\end{center}

\begin{center}\small\color{gray}对话 Prompt 的三个层次\end{center}

系统层的 Prompt 通常写在项目记忆文件中（如 CLAUDE.md），不需要每次对话重复。任务层是你在对话中最常使用的——它应该包含\textbf{“做什么”“用什么技术”“有什么约束”}三个要素。反馈层是最考验功力的——好的反馈能让 AI 精准修正，坏的反馈只会让它重新猜测。

\begin{tcbinfo}{小贴士：反馈三板斧}
一个实用的反馈模板是“指出问题 + 说明期望 + 提供上下文”。比如：“第 25 行的错误处理只用了 console.log（指出问题），生产环境应该用统一的错误处理器并返回用户友好提示（说明期望），项目里有一个 handleError 函数可以用（提供上下文）”。这种反馈比“改一下错误处理”有效十倍。

\end{tcbinfo}

\booksection{二、调试对话：如何描述错误让 AI 快速定位}

调试是 AI 编程中最频繁的场景之一。当你遇到 bug 时，如何向 AI 描述问题，直接决定了它能否快速帮你定位和修复。以下是调试对话的四个要点：

\begin{enumerate}
  \item \textbf{给 AI"看"原始信息，而不是你的转述}。把错误日志、异常堆栈、相关代码直接 @ 引用给 AI，而不是用"好像有个什么报错"来描述。AI 看到原始信息比看到你的理解更准确。
  \item \textbf{描述"期望行为"和"实际行为"的差异}。"我期望点击按钮后跳转到首页（期望），但实际点击后页面没有反应（实际）"——这种描述比"按钮坏了"有用得多。
  \item \textbf{提供复现步骤}。"先输入用户名，然后点击登录按钮，就会出现这个错误"——复现步骤让 AI 能模拟你的操作路径，更容易定位问题。
  \item \textbf{说清楚你已经尝试了什么}。"我已经检查了网络请求，返回 200，数据正常"——这能帮 AI 排除已验证的方向，不浪费时间重复排查。
\end{enumerate}

差的调试描述

\begin{lstlisting}
按钮点击没反应，帮我看看
\end{lstlisting}

好的调试描述

\begin{lstlisting}
点击登录按钮后页面无反应。
期望行为：点击后跳转到 /home 页面
实际行为：点击后无任何反应，控制台无报错
复现步骤：
1. 在登录页输入用户名admin和密码123456
2. 点击"登录"按钮
3. 页面无反应
已排查：
- 网络请求正常，返回200
- onClick事件已绑定
- 检查了LoginButton组件的代码
相关文件：@src/components/LoginButton.tsx
@src/pages/Login.tsx
\end{lstlisting}

\booksection{三、迭代策略：小步快跑 vs 大步重构}

AI 编程中的迭代有两种基本策略，适用于不同场景。理解它们的区别，能帮你在不同情况下做出正确选择。

\textbf{“小步快跑”}策略适用于功能开发和 bug 修复。它的核心是：每次让 AI 只做一个小的修改，验证通过后再做下一个。优势是每一步都可验证、可回退，风险可控；劣势是速度较慢，需要频繁的对话交互。适合：新功能开发、不确定最优方案时的探索、对 AI 产出不太放心时的谨慎迭代。

\textbf{“大步重构”}策略适用于明确知道目标的大规模修改。它的核心是：一次性描述完整的目标，让 AI 自主规划和执行整个修改。优势是速度快，AI 能看到全局图景做出更优的架构决策；劣势是风险高，如果 AI 理解偏差，修正成本大。适合：明确的架构升级、大规模代码迁移、有充分测试保障的重构。

\begin{center}
\begin{tabularx}{\linewidth}{|X|X|X|}
\hline
\textbf{维度} & \textbf{小步快跑} & \textbf{大步重构} \tabularnewline
\hline
每次修改量 & 一个功能点或一个bug & 一个模块或一个架构层 \tabularnewline
验证频率 & 每次修改后立即验证 & 全部修改完成后统一验证 \tabularnewline
对话轮次 & 多（5-10轮/功能） & 少（1-3轮/模块） \tabularnewline
风险 & 低，每步可回退 & 高，偏差修正成本大 \tabularnewline
适用场景 & 新功能、探索性开发 & 明确目标的重构、迁移 \tabularnewline
前提条件 & 无特殊要求 & 需要有测试覆盖 \tabularnewline
\hline
\end{tabularx}
\end{center}

\begin{center}\small\color{gray}两种迭代策略的对比\end{center}

实践中，两种策略常常组合使用。比如重构一个模块时，先用“大步重构”让 AI 给出整体方案，审查方案后，再用“小步快跑”逐步执行。这样既保证了全局视角，又控制了执行风险。

\booksection{四、多轮对话的上下文管理}

随着对话轮次增加，你需要主动管理对话上下文，避免“AI 忘记了早期信息”的问题。以下是几个实用的上下文管理技巧。

\textbf{第一，“锚定”关键信息。}在对话开始时，把关键约束写清楚，并在后续对话中适时重申。比如“记住，我们的数据存在 localStorage，不能用数据库”——在相关的修改请求中重申这个约束，防止 AI 遗忘。

\textbf{第二，善用“总结-续接”技巧。}当对话变得很长时，可以让 AI“总结一下目前的进展和待办事项”，然后开一个新对话，把这个总结作为起点。这样既保留了关键上下文，又避免了长对话的退化。

\textbf{第三，“一个对话一个任务”。}这是最简单也最有效的上下文管理策略。每个对话只处理一个任务或一个功能点，完成后开新对话处理下一个。这天然限制了对话长度，保持上下文清晰。

\textbf{第四，利用工具的上下文管理功能。}大多数 AI 编程工具都有“清除上下文”或“新会话”按钮——当你感觉 AI 的回答质量下降时，不要犹豫，开一个新会话。丢失一些对话历史比继续在退化的上下文中挣扎要好得多。

\begin{tcbpractice}{来动手：给你的 AI 当一次"病人"}
故意在浏览器里做一个小页面，制造一个 bug（比如按钮没反应）。然后按本章“好的调试描述”的格式，向 AI 完整描述一次：原始信息、期望 vs 实际、复现步骤、已排查项。感受一下 AI 定位问题的速度差异——你会理解为什么“描述问题”本身是门手艺。

\end{tcbpractice}

\begin{tcbsummary}{本章小结}
\begin{itemize}
  \item Prompt 三层：系统层（写在记忆文件）、任务层（做什么+技术+约束）、反馈层（指出问题+期望+上下文）
  \item 调试四要点：给原始信息、期望vs实际、复现步骤、已排查项
  \item 小步快跑：低风险高轮次，适合新功能；大步重构：高风险低轮次，需要测试保障
  \item 上下文管理四招：锚定关键信息、总结-续接、一个对话一个任务、果断开新会话
\end{itemize}

\end{tcbsummary}
